% =====================================================================
%  Deliverable 3 — Generative Artificial Intelligence (580694)
%  Universidad de Concepción · 31 de octubre de 2026
%  Joaquín Sepúlveda · Renato Saavedra · Lucas Saldivia
%
%  UNA PÁGINA, VERTICAL, compilado desde LaTeX. Overleaf → pdfLaTeX.
%  Las cifras de los 150 casos salen de scripts/analisis_d3.py y las de los
%  180 de scripts/analisis_replica_v4.py, sobre resultados/.
% =====================================================================
\documentclass[9pt,a4paper]{extarticle}

\usepackage[T1]{fontenc}
\usepackage[utf8]{inputenc}
\usepackage[spanish,es-noquoting]{babel}
% babel-spanish vuelve activa la comilla doble: "c se convierte en ç y rompe
% el JSON de ejemplo. Esto la desactiva.
\shorthandoff{"}
\usepackage{libertine}
\usepackage[scaled=0.80]{beramono}
\usepackage{microtype}
\usepackage[margin=8mm]{geometry}
\usepackage{multicol}
\usepackage{xcolor}
\usepackage{booktabs}
\usepackage{tikz}
\usepackage{array}
\usepackage{colortbl}
\usetikzlibrary{arrows.meta,positioning,fit,calc}
\usepackage[hidelinks]{hyperref}
\urlstyle{tt}

\definecolor{tinta}{HTML}{12191A}
\definecolor{acento}{HTML}{B8501C}
\definecolor{verde}{HTML}{146B5F}
\definecolor{ambar}{HTML}{C08A24}
\definecolor{gris}{HTML}{5A6A6B}
\definecolor{fondo}{HTML}{EDF1F0}
\definecolor{linea}{HTML}{C8D4D1}

\color{tinta}
\pagestyle{empty}
\setlength{\columnsep}{5mm}
\setlength{\parindent}{0pt}
\setlength{\parskip}{0.7mm plus 1mm}
\renewcommand{\arraystretch}{0.93}
\setlength{\tabcolsep}{3.2pt}

\newcommand{\seccion}[1]{\vspace{0.9mm}{\bfseries\footnotesize\color{acento}\MakeUppercase{#1}}\par\nobreak\vspace{0.4mm}\nobreak}
\newcommand{\q}{\textquotedbl}
\newcommand{\cod}[1]{{\ttfamily\small #1}}
\newcommand{\destaque}[1]{%
  \begingroup\setlength{\fboxsep}{2.0mm}%
  \colorbox{fondo}{\parbox{\dimexpr\columnwidth-4.0mm\relax}{#1}}%
  \endgroup\par}

\begin{document}
\raggedright

% ---------------------------------------------------------------- cabecera
{\Large\bfseries Tres agentes sobre el mismo modelo de 3,8B suben la validación de inscripción de ramos de 64,0\,\% a 88,0\,\%, y de 65,6\,\% a 91,1\,\% en 180 casos nuevos}\\[0.8mm]
{\footnotesize\color{gris}%
\textbf{Deliverable 3} · Generative Artificial Intelligence (580694) · Universidad de Concepción · 31 de octubre de 2026\\
\textbf{Equipo:} Joaquín Sepúlveda · Renato Saavedra · Lucas Saldivia \quad\textbullet\quad
\textbf{Modelo:} \cod{microsoft/Phi-3.5-mini-instruct} (revisión \cod{2fe1924}), 3,8B, NF4, T4 de Colab, en los tres agentes\\
\textbf{Repositorio:} \href{https://github.com/Luquitas02/genai-inscripcion-ramos}{\cod{github.com/Luquitas02/genai-inscripcion-ramos}} \quad\textbullet\quad
\textbf{Video:} {\small\url{PEGAR-LINK-AQUI}}}

\vspace{1.5mm}
{\color{linea}\rule{\textwidth}{0.4pt}}
\vspace{1mm}

\begin{multicols}{2}

% ================================================================ COLUMNA 1
\seccion{Tarea, falla y continuidad}

{\footnotesize
La tarea y el criterio son los del D1 y el D2. Dado el historial de un estudiante y una
pregunta en prosa, el sistema decide si puede inscribir el ramo y cita la regla. Un caso es
correcto si acierta decisión y regla. En el sistema del D2 el modelo solo leía la pregunta y
el código aplicaba el reglamento. Con el ramo y el período correctos acertaba 28 de 28, así que
todo su error venía de esa lectura. El D3 ataca esa falla con Phi, como se comprometió en el D2.

Separando los errores del D2 por mecanismo aparecen tres. \textbf{M1}, copiar el código: con el
nombre exacto del ramo en la pregunta, Phi acertaba 18 de 34, porque copiaba mal un código de
seis dígitos desde una lista de 61. \textbf{M2}, aislar la mención: si la pregunta nombra otro
ramo antes del consultado, acertaba 3 de 20. \textbf{M3}, el período, que fallaba cuando el
ramo de contexto se cursa \q ahora\q.\par}

\seccion{El grafo de agentes}

\begin{center}
\begin{tikzpicture}[
  font=\scriptsize,
  caja/.style={draw=linea,fill=white,rounded corners=1pt,inner sep=1.3mm,text width=33mm,align=center},
  agente/.style={caja,draw=acento,line width=0.6pt},
  herr/.style={caja,draw=verde,line width=0.6pt,text width=31mm},
  fl/.style={-{Stealth[length=1.5mm]},draw=gris},
  et/.style={font=\tiny,text=gris,align=left},
  node distance=5.2mm]

\node[caja,fill=fondo] (q) {pregunta en prosa};
\node[agente,below=of q] (a) {\textbf{A · Intérprete}\\ lee la pregunta, no ve catálogo};
\node[herr,below=7mm of a] (t1) {\textbf{T1 · buscar\_ramo}\\ código, sin modelo};
\node[agente,below=of t1] (b) {\textbf{B · Resolvedor}\\ elige entre 5 candidatos};
\node[agente,below=of b] (c) {\textbf{C · Verificador}\\ contrasta con la pregunta};
\node[herr,below=of c] (t2) {\textbf{T2 · motor de reglas}\\ el código del D2};
\node[caja,fill=fondo,below=of t2] (out) {\textbf{fin} · respuesta};
\node[caja,draw=gris,dashed,text width=12mm,right=9mm of a] (fin) {\textbf{fin} sin ramo};

\draw[fl] (q) -- node[right,et]{texto} (a);
\draw[fl] (a) -- node[right,et]{busca: consultado (otros y\\ período siguen a B, C y T2)} (t1);
\draw[fl] (t1) -- node[right,et]{5 candidatos\\ con puntaje} (b);
\draw[fl] (b) -- node[right,et]{elige: código} (c);
\draw[fl] (c) -- node[right,et]{acepta: código, período} (t2);
\draw[fl] (t2) -- node[right,et]{decisión y regla} (out);
\draw[fl,dashed] (a) -- node[above,et]{sin\_ramo} (fin);
% bucle del Resolvedor con T1
\draw[fl] (b.east) -- ++(4.5mm,0) |- node[pos=0.25,right,et]{busca de nuevo,\\ nueva consulta\\ (máx. 2)} (t1.east);
% tope: el caso sigue con la última elección
\draw[fl,dashed] (c.east) -- ++(4.5mm,0) |- node[pos=0.25,right,et]{tope (3.er rechazo\\ o 12 llamadas):\\ última elección} (t2.east);
% rechazo del Verificador al Intérprete
\draw[fl] (c.west) -- ++(-15mm,0) |- node[pos=0.3,left,et,align=right]{rechaza:\\ problema\\ y motivo;\\ vuelve\\ máx. 2} (a.west);
\draw[fl] (b.west) -- ++(-3mm,0) |- node[pos=0.25,left,et,align=right]{ninguno:\\ motivo} ($(a.west)+(0,-1.2mm)$);
\end{tikzpicture}
\end{center}

{\footnotesize\color{gris} Borde naranjo, agentes, cada uno con su archivo en \cod{prompts/}.
Borde verde, herramientas. Las vueltas al Intérprete, por rechazo o por \cod{ninguno}, suman
máx. 2. Tope de 12 llamadas por caso. Al llegar a un tope se usa la última elección.\par}

\vspace{0.6mm}
{\footnotesize Una traza real (nuevos, caso 49), con las salidas de cada agente:\par}
\vspace{0.3mm}
{\setlength{\fboxsep}{1.4mm}\colorbox{fondo}{\parbox{\dimexpr\columnwidth-2.8mm\relax}{\ttfamily\fontsize{6.6}{8}\selectfont\raggedright
\textrm{\footnotesize\q me falta programación pero la estoy cursando, ¿puedo tomar opti I el que viene?\q}\\[0.4mm]
A$\to$T1 \{\q{}consultado\q{}: \q{}opti I\q{}, \q{}otros\q{}: [\q{}programación\q{}], \q{}frase\_periodo\q{}: \q{}el que viene\q{}, \q{}periodo\q{}: \q{}proximo\q{}\}\\
T1$\to$B 580315 Optimización I 0,60 · 201 Complementaria 1 0,55 · ...\\
B$\to$T1 \{\q{}accion\q{}: \q{}buscar\q{}, \q{}consulta\q{}: \q{}Optimización I\q{}\}\\
T1$\to$B 580315 Optimización I 1,50 · 580325 Optimización II 0,93 · ...\\
B$\to$C \{\q{}accion\q{}: \q{}elegir\q{}, \q{}codigo\q{}: \q{}580315\q{}\}\\
C$\to$T2 \{\q{}cita\_ramo\q{}: \q{}Optimización I\q{}, \q{}accion\q{}: \q{}aceptar\q{}\}\\
T2 condicional · R-DEPENDE-APROBACION (correcta)}}\par}

\seccion{Los agentes}

{\footnotesize
\begin{tabular}{@{}p{12mm}p{37mm}p{38mm}@{}}
\toprule
agente & recibe y devuelve & decide \\
\midrule
Intérprete & la pregunta $\to$ ramo consultado con las palabras del estudiante, los otros ramos, la frase que dice cuándo y el período & entregar o \cod{sin\_ramo}. Ante un rechazo, revisa su respuesta. No usa herramientas \\
Resolvedor & la mención y 5 candidatos de T1 $\to$ un código & elegir, llamar a T1 con otra consulta o \cod{ninguno} \\
Verificador & la pregunta y la elección $\to$ citas de la pregunta y su veredicto & aceptar y llamar a T2, o rechazar con motivo \\
\bottomrule
\end{tabular}\par}

\vspace{0.6mm}
{\footnotesize
El trabajo se divide por mecanismo y por contexto. El Intérprete no ve el catálogo, así que no
puede caer en M1. El Resolvedor nunca ve más de cinco candidatos. El Verificador es el único que
ve juntas la pregunta y la elección. Las reglas siguen en código porque el D2 midió que el modelo
las aplica al 31,7\,\% con la ficha perfecta y el código al 100\,\%.\par}

% ================================================================ COLUMNA 2
\columnbreak

\seccion{Resultados}

{\footnotesize
\begin{tabular}{@{}lrrrrrr@{}}
\toprule
& \multicolumn{3}{c}{corrida del 20 oct} & \multicolumn{3}{c}{réplica, casos sellados} \\
\cmidrule(lr){2-4}\cmidrule(l){5-7}
acierto conjunto & plantilla & redact. & \textbf{150} & plantilla & redact. & \textbf{180} \\
\midrule
baseline few-shot (D1)   &  28 &  7 &  35 &  30 & 13 &  43 \\
sistema del D2           &  72 & 24 &  96 &  76 & 42 & 118 \\
regla fija, sin modelo   &  83 &  7 &  90 &  77 & 25 & 102 \\
\rowcolor{fondo}
\textbf{tres agentes}    & \textbf{103} & \textbf{29} & \textbf{132} & \textbf{109} & \textbf{55} & \textbf{164} \\
\quad con ramo y período bien & 97 & 24 & 121 & 103 & 55 & 158 \\
\bottomrule
\end{tabular}\par}

\vspace{1mm}
\destaque{\footnotesize En los 150 casos los agentes suben de 64,0\,\% a 88,0\,\%. Donde difieren con el D2 ganan 39 y pierden 3 (McNemar exacto, $p<0{,}001$), pero en los 30 redactados la diferencia no era significativa (6 a 1, $p=0{,}125$). Para comprobarlo se corrió el mismo sistema, sin cambios y con un protocolo subido antes, sobre 180 casos sellados que nadie había usado, 60 de ellos redactados a ciegas. Ahí ganan 50 a 4 ($p<0{,}001$), y 14 a 1 en las redactadas.}

\vspace{0.6mm}
{\footnotesize
La regla fija toma el texto que sigue a \q inscribir\q, \q tomar\q\ o \q meter\q\ y lo busca con T1.
Iguala o supera al D2 en las preguntas de plantilla y cae a 7 de 30 y 25 de 60 en las redactadas.
Los agentes no dependen de la plantilla. Exigiendo también ramo y período correctos, el D2 acierta
75 y 94. Por mecanismo, en originales y nuevos sin ambiguos:\par}

\vspace{0.3mm}
{\footnotesize
\begin{tabular}{@{}lrrr@{}}
\toprule
acierta el ramo, cuando la pregunta... & casos & D2 & agentes \\
\midrule
trae el nombre exacto (M1)            & 69 & 44 & \textbf{67} \\
lo abrevia o dice coloquial           & 36 & 19 & \textbf{33} \\
\quad de ambos, nombra otro ramo antes (M2) & 35 &  7 & \textbf{34} \\
\midrule
acierta el período (M3)               & 105 & 90 & \textbf{98} \\
\bottomrule
\end{tabular}\par}

\seccion{Técnicas y ablación}

{\footnotesize
\begin{tabular}{@{}llllr@{}}
\toprule
técnica (Lecture 5) & agente & ataca & sin ella & $\Delta$ \\
\midrule
uso de herramienta (ReAct)   & B & M1      & 75  & $-57$ \\
few-shot, 4 ejemplos         & A & M2      & 123 & $-9$ \\
descomposición en menciones$^{*}$ & A & M2 & 136 & $+4$ \\
Self-Refine con citas        & C & M2, M3  & 133 & $+1$ \\
B reemplazado por el top-1   & B & M1      & 133 & $+1$ \\
\bottomrule
\end{tabular}\par}

\vspace{0.6mm}
{\footnotesize
Cada fila es el sistema completo (132) sin una pieza, sobre los mismos 150 casos. Las ablaciones
no se repitieron en la réplica. ($^{*}$) Esta fila quita dos instrucciones a la vez, la lista de
los otros ramos y la regla de tomar el ramo que acompaña al verbo. Todos los agentes usan además
un rol y salida JSON. El Verificador lleva 2 ejemplos, y el Intérprete y el Verificador copian
la frase de la pregunta antes de decidir (\cod{frase\_periodo}, \cod{cita\_ramo}). Nada de eso se
ablacionó.

La mejora la explican dos piezas. Sin T1 el modelo vuelve a copiar códigos y el sistema cae a
75 (McNemar $p<0{,}001$). Sin los ejemplos cae a 123 ($p=0{,}004$) porque M2 baja de 34 a 26.
El período, en cambio, sube sin ejemplos (141 contra 137).

El Resolvedor y el Verificador no aportan, y las trazas muestran por qué. El Resolvedor se apartó
del primer candidato de T1 en 8 de los 150 casos y en 8 de los 180. Acertó en 1 de esos 16, y
en 8 el primer candidato era el correcto. Un rechazo del Verificador vuelve al Intérprete. Si su
mención estaba bien, la repite, T1 devuelve los mismos candidatos y el Resolvedor elige lo mismo.
Por eso 4 de los 5 rechazos de cada corrida terminan en tope. En originales 59 el Verificador
rechazó con razón tres veces Optimización II y el tope aceptó esa elección.\par}

\seccion{Caso de falla}

{\footnotesize
Elegido por la regla del protocolo: el primer redactado que falla. \q ...ando cursando
administración y quería tomar marketing el otro semestre, en marzo\q. El Intérprete copió
\q en marzo\q\ y respondió \q proximo\q, que era correcto. El Verificador rechazó con el motivo
\q el otro semestre sería el semestre actual\q. El Intérprete cambió a \q actual\q, el
Verificador aceptó y T2 respondió \q no\q\ por el prerrequisito en curso. La respuesta correcta
era \q condicional\q. Un modelo de 3,8B que revisa a otro puede deshacer un acierto, y es el
único caso de los 150 en que el Verificador cambió el resultado.\par}

\seccion{Cómo se midió}

{\footnotesize
Los prompts se ajustaron en 52 casos de desarrollo (de 40 a 44 en cuatro versiones) y en 21
preguntas de desarrollo redactadas (20 de 21). Las reglas para aceptar cada versión y los
protocolos de las dos corridas se subieron al repositorio antes de correr. Las preguntas redactadas las escribió un modelo de lenguaje sin acceso
al sistema y el equipo las leyó. En las 30 primeras, Claude, que conocía el sistema, corrigió la
redacción de 3. Cada mensaje entre agentes queda en \cod{trazas/} con el hash del prompt que lo
produjo.

Después de la corrida del 20 de octubre probamos una versión 5, con un Verificador que corrige
el período y devuelve el ramo al Resolvedor. En 223 casos de desarrollo sus dos variantes sacaron
186 y 194 contra 196 de esta versión, y no se aceptó.\par}

\seccion{Límites declarados}

{\footnotesize
\textbf{Ambiguos.} En 15 de los 120 casos de plantilla la pregunta nombra el ramo con una forma
que sirve para varios, como \q gestión\q\ o \q la I\q. Ningún lector de la pregunta los resuelve.
\textbf{T1 y los 150.} El puntaje de T1 se ajustó mirando las menciones de originales y nuevos.
Las redactadas y los 180 sellados no los vio.
\textbf{Casos parecidos.} 9 casos de desarrollo comparten historial y consulta con 7 de los 150,
y 38 de los 180 repiten el problema o la pregunta de un caso anterior. Sin esos 38, los agentes
ganan 34 a 4 (cálculo exploratorio).
\textbf{Costo.} 3,2 llamadas, 2.797 tokens y 7,2\,s por caso, contra 1.155 tokens y 1,7\,s del D2.\par}

\seccion{Lo que muestra el video}

{\footnotesize
Tres casos de los 150 elegidos por la regla del protocolo, corridos en vivo con
\cod{demo\_d3\_colab.ipynb}, que compara cada mensaje con la traza guardada. Uno pasa por los
tres agentes sin bucles. En otro el Verificador rechaza el período tres veces y el tope cierra
con la respuesta correcta. El tercero es el caso de falla.\par}

\end{multicols}
\end{document}
